% =======================================================================================
% Appendice: riproducibilità
% =======================================================================================
\chapter{Riproducibilità dei risultati}\label{app:reproducibility}
\begingroup\globaldefs=1 \pagestyle{appendices}\endgroup

\noindent Questa appendice documenta come i risultati presentati nel
\autoref{chap:results} possano essere riprodotti e verificati.

% =======================================================================================
\section{Struttura del materiale}\label{sec:appendix-layout}

Il materiale è organizzato nelle seguenti cartelle.

\begin{description}[leftmargin=0em]
    \item[\texttt{RENUVER/}] l'archivio di partenza: eseguibile, dataset, file di regole,
    materiale di ground truth e file di risultato delle esecuzioni originali. L'eseguibile non è
    mai stato modificato; la verifica di integrità è effettuata per checksum. L'archivio include
    un \texttt{README} che documenta la struttura delle cartelle e i parametri della riga di
    comando: le discrepanze fra quel documento e il comportamento osservato del codice sono
    discusse nella \autoref{sec:architecture} e nel \autoref{chap:defects}.
    \item[\texttt{pruning/}] gli strumenti sviluppati: il pruner, il sistema di valutazione, gli
    script di verifica e quelli di orchestrazione delle campagne.
    \item[\texttt{results/}] i file di risultato di tutte le campagne, in formato JSON, con una
    voce per ogni esecuzione e i relativi indicatori; in \texttt{results/analisi/} i dati delle
    prove che non appartengono a una campagna.
    \item[\texttt{analisi/}] l'analisi del sistema: disassemblaggio commentato, modello dei dati,
    inventario del codice non raggiungibile.
    \item[\texttt{report.md}] la documentazione estesa del lavoro, con la derivazione completa
    delle condizioni teoriche e il dettaglio di tutte le campagne.
\end{description}

% =======================================================================================
\section{Verifica automatica dei risultati}\label{sec:appendix-audit}

I numeri asseriti nella documentazione sono ricalcolati dai file di risultato da una procedura
automatica di verifica, che li confronta con i valori asseriti e segnala ogni divergenza. La
procedura copre \num{67} asserzioni, fra cui le condizioni teoriche derivate dal bytecode, le
metriche aggregate di ciascuna campagna e i risultati dei test statistici.

La procedura è stata impiegata durante la stesura: ha individuato e consentito di correggere tre
valori non aggiornati, fra cui un risultato il cui segno si invertiva estendendo il campione da
\num{14} a \num{15} configurazioni, e che era stato inizialmente riportato con il segno errato.
La circostanza motiva l'adozione della procedura: un risultato che cambia segno aggiungendo una
configurazione su quindici non è robusto, e non deve essere presentato come se lo fosse.

Due procedure ulteriori completano la verifica dei numeri citati nella tesi, ed entrambe operano
sui soli file di risultato, senza rieseguire il sistema.

\begin{sloppypar}
\begin{itemize}[leftmargin=1.4em]
    \item la prima, \texttt{stats\_effects.py}, ricalcola i test della famiglia primaria con la
    medesima ricetta di appaiamento impiegata dalla campagna --- chiave (dataset, soglia,
    ripetizione), baseline e variante entrambe concluse con esito regolare --- e produce, oltre ai
    valori di $p$ già pubblicati, la dimensione dell'effetto, l'intervallo di confidenza bootstrap
    della differenza mediana e il valore di $p$ corretto con la procedura di Holm
    (\autoref{tab:statistics});
    \item la seconda, \texttt{make\_appendix\_tables.py}, genera le tabelle complete
    dell'\autoref{app:results} dai file di risultato; i totali che esse riportano coincidono con
    i valori discussi nel \autoref{chap:results}, e la generazione è quindi anche un controllo di
    coerenza fra i numeri citati e i dati grezzi.
\end{itemize}
\end{sloppypar}

Le due procedure sono deterministiche: a parità di file di ingresso producono lo stesso output, e
il ricampionamento impiegato per gli intervalli di confidenza usa un seme fissato. I dati
dell'esperimento di ordinamento, che non appartengono a una campagna ma a un controllo dedicato,
sono archiviati con gli altri in \texttt{results/analisi/}.

% =======================================================================================
\section{Verifica di equivalenza della ricompilazione}\label{sec:instrumentation}

La caratterizzazione del costo computazionale (\autoref{sec:cost}) ha richiesto di ricompilare il
sistema, distribuito senza sorgenti. La procedura adottata è la seguente.

\begin{enumerate}[leftmargin=1.6em]
    \item \textbf{Decompilazione} delle classi applicative dell'archivio.
    \item \textbf{Ricompilazione} per la versione \num{52} del formato \texttt{class}, contro le
    medesime librerie incorporate nell'archivio originale.
    \item \textbf{Ricostruzione} dell'archivio sostituendo \emph{soltanto} le classi applicative
    dentro una copia dell'originale: le restanti classi di libreria restano quelle di partenza.
    \item \textbf{Verifica di equivalenza} eseguendo l'archivio originale e quello ricostruito
    sulla medesima configurazione, in cartelle isolate, e confrontando gli output.
\end{enumerate}

L'esito della verifica è riportato nella \autoref{tab:equivalence}: i file di output coincidono
byte per byte, e le uniche differenze riguardano i tempi di esecuzione registrati, che variano
fisiologicamente fra esecuzioni.

\begin{table}[htbp]
\centering
\caption{Verifica di equivalenza della ricompilazione. Configurazione
\texttt{bridges\_70\_1}, soglia 15.}
\label{tab:equivalence}
\footnotesize
\begin{tabular}{lll}
\toprule
file & dimensione & esito del confronto \\
\midrule
\texttt{Populated} & \SI{1377}{\byte}  & \textbf{identico} (checksum \texttt{682428f09943}) \\
\texttt{Results}   & \SI{11522}{\byte} & \textbf{identico} (checksum \texttt{8eca4ca48f2a}) \\
\texttt{Logs}      & ---               & differisce \emph{solo} nelle righe di tempo \\
\bottomrule
\end{tabular}
\end{table}

Solo dopo questa verifica il sistema è stato strumentato, aggiungendo contatori e misure di tempo
mediante \emph{wrapper} sui metodi di interesse anziché modificandone i corpi, in modo che il
flusso di controllo restasse quello originale.

% =======================================================================================
\section{Comandi principali}\label{sec:appendix-commands}

Le campagne sono orchestrate da script che invocano il sistema di valutazione. A titolo
esemplificativo, la valutazione di una strategia su una famiglia di dataset richiede:

\begin{quote}
\texttt{python3 pruning/benchmark.py \textbackslash}\\
\texttt{\ \ --datasets <elenco> --runs <elenco> --thresholds <elenco> \textbackslash}\\
\texttt{\ \ --variants "<nome>=<strategia>:<politica>" \textbackslash}\\
\texttt{\ \ --out results/<nome>.json}
\end{quote}

\noindent Il sistema di valutazione provvede a eseguire la configurazione di riferimento, a
produrre il file di regole potato per ciascuna variante, a eseguire il sistema su entrambi e a
calcolare le metriche in forma appaiata. I risultati sono memorizzati incrementalmente, cosicché
una campagna interrotta può essere ripresa.

% =======================================================================================
\section{Ambiente di esecuzione}\label{sec:appendix-env}

Le campagne sono state condotte su una macchina a otto core, con esecuzioni strettamente
sequenziali. La sequenzialità non è una scelta di comodo: la cartella di lavoro del sistema è
condivisa fra esecuzioni, ed esecuzioni simultanee sovrascrivono reciprocamente gli output
producendo risultati errati senza segnalare alcun errore. La circostanza è stata verificata
sperimentalmente.

Le versioni delle librerie impiegate per la ricompilazione sono fissate e documentate, poiché
versioni diverse possono produrre differenze nel comportamento della generazione dei numeri
casuali e quindi nei risultati.
